\SourceLeaf{206}The account can therefore be framed as follows. The equation
\[
\left.\begin{array}{c} S \\ \Sigma \end{array}\right\} = M\pm h+i+u
\]
shows that equal positions of the bubble, with $i$ unchanged, indicate equal values of $u$. Now the values of $u$ in I and II are opposite. If [the bubble] therefore has the same position in II as in I, then $u'=u''=0$.

In the opposite case we conclude that $u$ was not $=0$ and has now taken the opposite of its previous value. If we therefore return $u$ to its previous value by steadily turning screw 7, it passes through 0 in the process, and turning the screw halfway back brings it to 0. The line 1\,2 or 2\,1 or 3\,4 is then horizontal, though the bubble itself has likewise moved halfway back toward its 2$^{\mathrm{nd}}$ position. We now change $i$ by turning screw 3 until $\left.\substack{S\\\Sigma}\right\} = M\pm h$.

\begin{center}
\textsc{Remarks.}
\end{center}

The first of the preceding notes is found on the rear flyleaf of \emph{Mathematische Abhandlungen} by Jacob Wilhelm Heinrich \textsc{Lehmann}, Zerbst 1829, and presumably dates from that year. Note [2] is on a separate leaf. The date of its composition cannot be established exactly, though it may be presumed to fall between 1820 and 1830. \textsc{Gauss} recorded note [3] in a notebook, most likely in March 1832, since it refers to \textsc{Gauss}'s letter to \textsc{Bolyai} of 6 March 1832. Note [4] is based on an oral communication from \textsc{Gauss} to \textsc{Lübsen}, who attended his lectures in 1830 (see \textsc{Lübsen}, \emph{Ausf\"uhrliches Lehrbuch der Analysis}, Hamburg 1853, p. 171). Finally, note [5] is written on an undated slip of paper. The handwriting points to a later period of \textsc{Gauss}'s life, roughly 1840--1850. \textsc{Gauss} probably wrote the note while preparing a lecture on practical astronomy. Its content alone determined its placement here.

That \textsc{Gauss} concerned himself with the definition of the plane very early is shown by the passage already printed on p. 162 of this volume from his diary entry of 28 July 1797.

The two figures in note [1] were added to the text.

\hfill \textsc{Stäckel.}

